\documentclass[%
	12pt,%11pt, 10pt%
	a4paper,%a5paper, letterpaper, legalpaper, executivepaper
%	onelinecaption,% noonelinecaption,%
	oneside,%twoside,%
%	onecolumn,%twocolumn,%
% openany, openright,%
% cleardoublestandard,% cleardoubleplain, %cleardoubleempty
% notitlepage,%titlepage,%
	listof=totoc, index=totoc, bibliography=totoc, %bibtotocnumbered, liststotocnumbered
% tocindent,% tocleft,%
% listsindent,% listsleft%
	parskip=half,% parindent, parskip, parskip*, parskip+, parskip, halfparskip*, halfparskip+, halfparskip,
	chapterprefix=false,%chapterprefix,%
	appendixprefix, %noappendixprefix,%
%	headsepline,headnosepline%
%	footsepline,footnosepline%
% bigheadings, normalheadings,
	headings=small,%
%%Formatierungsparameter
%	abstractoff,% abstracton%
% pointednumbers,% pointlessnumbers%
% leqno, fleqn%
%	tablecaptionbelow,%
	captions=tableheading,%
% origlongtable,%
% openbib,%
% final, draft%
]{scrreprt}
%]{scrbook}
%]{scrartcl}
%]{scrlttr2}
%
%]{article}
%]{report}
%]{letter}
%]{book}
%]{slides}

\usepackage{bibgerm}					% deutsche Bibliography
\usepackage[ngerman]{babel}		% Sprachen
\usepackage[utf8]{inputenc} % Eingabe von ü,ö,ä,ß erlaubt
\usepackage[T1]{fontenc}			% Umlauten koennen richtig getrennt werden

\usepackage{lmodern}

%======================================================================
%   Darstellung des Glossars und des Indexes
%======================================================================
\usepackage{makeidx}   %Änderung
\usepackage{booktabs}
%\usepackage{caption}							% mehrzeilige Captions ausrichten
%\captionsetup{format=hang,margin=5pt,labelfont=bf,position=bottom,aboveskip=1em,belowskip=1em}
%%\setcapindent{-1em}							 	% Zeilenumbruch bei Bildbeschreibungen.
%\setlength{\belowcaptionskip}{1em}

%\usepackage{multicol}            	% mehrere Spalten

%\usepackage{setspace}            % Zeilenabstand einstellbar
%\onehalfspacing                  % eineinhalbzeilig einstellen

%\usepackage{theorem}
%\theoremstyle{break}
\newtheorem{definition}{Definition}
%\usepackage{varioref}%

\usepackage{float}   %Änderung

%\usepackage{array}
%\usepackage{tabularx}
%\usepackage{longtable}

%-------------------------------Farbdefinitionen------------------------------------------------------
\usepackage{color}    %Änderung
\definecolor{sun1}{rgb}{0.2,0.2,0.4} %dunkelblau
\definecolor{sun2}{rgb}{0.4,0.4,0.6} %mittel
\definecolor{sun3}{rgb}{0.6,0.6,0.8} %mittel
\definecolor{sun4}{rgb}{0.8,0.8,1.0} %hellblau
\definecolor{sun5}{rgb}{0.9,0.9,1.0} %hellblau
\definecolor{black}{rgb}{0,0,0} %hellblau

%-------------------------------Kopf- und Fußzeile----------------------------------------------------
%\usepackage{scrpage2}
%\pagestyle{scrheadings}
%\pagestyle{scrplain}

%\ihead[scrplain-innen]{scrheadings-innen}
%\chead[scrplain-zentriert]{scrheadings-zentriert}
%\ohead[scrplain-außen]{scrheadings-außen}
%\ifoot[scrplain-innen]{scrheadings-innen}
%\cfoot[scrplain-zentriert]{scrheadings-zentriert}
%\ofoot[scrplain-außen]{scrheadings-außen}
%
%\usepackage{fancyhdr}
%\pagestyle{fancy}
%\addtolength{\headheight}{\baselineskip}

%\fancyhf{}   % löscht alle Felder
%
%\fancyhead[RO]{\nouppercase{\sffamily{\rightmark}}}%Kopfzeile rechts bzw. außen
%\fancyhead[LE]{\nouppercase{\sffamily{\leftmark}}}
%\renewcommand{\headrule}{\vbox to 0pt{\hbox to \headwidth{\dotfill}\vss}} %gepunktete Linie
%
%\fancyfoot[CO,CE]{\color{sun1}\sffamily{\thepage}}					%Fußzeile rechts bzw. außen
%\renewcommand{\footrule}{\vbox to 0pt{\hbox to \headwidth{\dotfill}\vss}} %gepunktete Linie
%
%%extra Formatierungen der Kopf- und Fußzeilen für Kapitelseiten
%\fancypagestyle{plain}{
%	\fancyhf{}   % löscht alle Felder
%
%	\fancyhead[LE,RO]{\nouppercase{\sffamily{\leftmark}}}%Kopfzeile rechts bzw. außen
%	\renewcommand{\headrule}{\vbox to 0pt{\hbox to \headwidth{\dotfill}\vss}} %gepunktete Linie
%
%	\fancyfoot[CO,CE]{\color{sun1}\sffamily{\thepage}}					%Fußzeile rechts bzw. außen
%	\renewcommand{\footrule}{\vbox to 0pt{\hbox to \headwidth{\dotfill}\vss}} %gepunktete Linie
%}

%======================================================================
%	Erkennen des Bearbeitungsmodus für weitere Paketoptionen
%======================================================================
\RequirePackage{ifthen}

\newif\ifhtm
\ifx \HCode\undefined \htmfalse \else	\htmtrue \fi

\newif\ifpdf
\ifx \pdfoutput\undefined
	\pdffalse 		%latex in DVI mode
\else
	\ifthenelse{\number\pdfoutput<1}%
 	{\pdffalse} 	%pdflatex in DVI mode
 	{\pdftrue} 		%pdflatex in PDF mode
\fi

%======================================================================
%		PDF <=> DVI <=> XHTML
%======================================================================
\ifhtm%
		% Bearbeitung mit tex4ht
		\usepackage{graphicx}%
		\usepackage[tex4ht]{hyperref}%
\else%
	\ifpdf%
		%Bearbeitung mit pdftex im PDF-Modus
		\usepackage[pdftex]{graphicx}%
		\usepackage[pdftex]{hyperref}%
		\pdfcompresslevel=0%						Damit wird die PDF-Quelldatei lesbar
		\pdfoptionpdfminorversion=6%		Bestimmt die PDF - Version der Ausgabe
	%\pdfadjustspacing=0%							0, 1 oder 2 Änderung nicht erkannt
%		\AtBeginDocument{\usepackage[tagged, flatstructure]{accessibility}}	% Makros in accessibility.sty
		\usepackage[tagged, highstructure]{accessibility}	% Makros in accessibility.sty
	\else%
		%Bearbeitung mit latex oder pdftex im DVI-Modus
		\usepackage{graphicx}%
		\usepackage{hyperref}%
		\newcommand{\alt}[1]{}%
		\newcommand{\thead}[1]{\textbf{#1}}
	\fi%
\fi%

% Einstellungen des Hyperref-Paketes
%\AtBeginDocument{%
	\hypersetup{%
%		pdftitle	= {\dctitle}, %
%		pdfsubject	= {\dcsubject, \dcdate}, %
%		pdfauthor	= {\dcauthorname~\dcauthorsurname, \dcauthoremail}, %
%		pdfkeywords	= {\dckeywords}, %
		pdfcreator	= {pdfTeX with Hyperref and Thumbpdf}, %
		pdfproducer	= {LaTeX, hyperref, thumbpdf}, %
		bookmarksnumbered,%
%		hypertexnames = true,%
%		plainpages = true,%
%		unicode = true,%
	}%
%}


%\usepackage[hyper=true, style=altlist, toc=true]{glossary}
%\usepackage{glossary}

%======================================================================
%\makeindex
%\makeglossary
%%makeindex -s minimaldokument.ist -t minimaldokument.glg -o minimaldokument.gls minimaldokument.glo
%======================================================================
\begin{document}

\setchapterpreamble{Ich bin die Präamble}
\chapter{listliketab}

Der große Oxmox riet ihr davon ab, da es dort wimmele von bösen Kommata, wilden Fragezeichen und hinterhältigen Semikoli, doch das Blindtextchen ließ sich nicht beirren. Es packte seine sieben Versalien, schob sich sein Initial in den Gürtel und machte sich auf den Weg. Als es die ersten Hügel des Kursivgebirges erklommen hatte, warf es einen letzten Blick zurück auf die Skyline seiner Heimatstadt Buchstabhausen, die Headline von Alphabetdorf und die Subline seiner eigenen Straße, der Zeilengasse. Wehmütig lief ihm eine rhetorische Frage über die Wange, dann setzte es seinen Weg fort.

\begin{table}[htbp]
	\centering
		\begin{tabular}{l|l l}
			\thead{11} & \thead{12} & \thead{13} \\ \hline
			11 & 12 & 13 \\
			21 & 22 & 23 \\
			31 & 32 & 33
		\end{tabular}
	\caption{meine ungetaggte Tabelle}
	\label{tab:meineUngetaggteTabelle}
\end{table}

Der große Oxmox riet ihr davon ab, da es dort wimmele von bösen Kommata, wilden Fragezeichen und hinterhältigen Semikoli, doch das Blindtextchen ließ sich nicht beirren. Es packte seine sieben Versalien, schob sich sein Initial in den Gürtel und machte sich auf den Weg. Als es die ersten Hügel des Kursivgebirges erklommen hatte, warf es einen letzten Blick zurück auf die Skyline seiner Heimatstadt Buchstabhausen, die Headline von Alphabetdorf und die Subline seiner eigenen Straße, der Zeilengasse. Wehmütig lief ihm eine rhetorische Frage über die Wange, dann setzte es seinen Weg fort.

\begin{figure}[htbp]
		\centering
		\alt{Ich bin das Logo der Technischen Universität Dresden}
		\includegraphics[scale=0.50]{images/tu_logo}%MCID 8
	\caption[tu]{tu-logo1}%MCID 9
	\label{fig:tulogo1}
\end{figure}

Der große Oxmox riet ihr davon ab, da es dort wimmele von bösen Kommata, wilden Fragezeichen und hinterhältigen Semikoli, doch das Blindtextchen ließ sich nicht beirren. Es packte seine sieben Versalien, schob sich sein Initial in den Gürtel und machte sich auf den Weg. Als es die ersten Hügel des Kursivgebirges erklommen hatte, warf es einen letzten Blick zurück auf die Skyline seiner Heimatstadt Buchstabhausen, die Headline von Alphabetdorf und die Subline seiner eigenen Straße, der Zeilengasse. Wehmütig lief ihm eine rhetorische Frage über die Wange, dann setzte es seinen Weg fort.

Der große Oxmox riet ihr davon ab, da es dort wimmele von bösen Kommata, wilden Fragezeichen und hinterhältigen Semikoli, doch das Blindtextchen ließ sich nicht beirren. Es packte seine sieben Versalien, schob sich sein Initial in den Gürtel und machte sich auf den Weg. Als es die ersten Hügel des Kursivgebirges erklommen hatte, warf es einen letzten Blick zurück auf die Skyline seiner Heimatstadt Buchstabhausen, die Headline von Alphabetdorf und die Subline seiner eigenen Straße, der Zeilengasse. Wehmütig lief ihm eine rhetorische Frage über die Wange, dann setzte es seinen Weg fort. Der große Oxmox riet ihr davon ab, da es dort wimmele von bösen Kommata, wilden Fragezeichen und hinterhältigen Semikoli, doch das Blindtextchen ließ sich nicht beirren. Es packte seine sieben Versalien, schob sich sein Initial in den Gürtel und machte sich auf den Weg. Als es die ersten Hügel des Kursivgebirges erklommen hatte, warf es einen letzten Blick zurück auf die Skyline seiner Heimatstadt Buchstabhausen, die Headline von Alphabetdorf und die Subline seiner eigenen Straße, der Zeilengasse. Wehmütig lief ihm eine rhetorische Frage über die Wange, dann setzte es seinen Weg fort. Der große Oxmox riet ihr davon ab, da es dort wimmele von bösen Kommata, wilden Fragezeichen und hinterhältigen Semikoli, doch das Blindtextchen ließ sich nicht beirren. Es packte seine sieben Versalien, schob sich sein Initial in den Gürtel und machte sich auf den Weg. Als es die ersten Hügel des Kursivgebirges erklommen hatte, warf es einen letzten Blick zurück auf die Skyline seiner Heimatstadt Buchstabhausen, die Headline von Alphabetdorf und die Subline seiner eigenen Straße, der Zeilengasse. Wehmütig lief ihm eine rhetorische Frage über die Wange, dann setzte es seinen Weg fort. Die Copy warnte das Blindtextchen, da, wo sie herkäme wäre sie zigmal umgeschrieben worden und alles, was von ihrem Ursprung noch übrig wäre, sei das Wort "`und"'  und das Blindtextchen solle umkehren und wieder in sein eigenes, sicheres Land zurückkehren.

Der große Oxmox riet ihr davon ab, da es dort wimmele von bösen Kommata, wilden Fragezeichen und hinterhältigen Semikoli, doch das Blindtextchen ließ sich nicht beirren. Es packte seine sieben Versalien, schob sich sein Initial in den Gürtel und machte sich auf den Weg. Als es die ersten Hügel des Kursivgebirges erklommen hatte, warf es einen letzten Blick zurück auf die Skyline seiner Heimatstadt Buchstabhausen, die Headline von Alphabetdorf und die Subline seiner eigenen Straße, der Zeilengasse. Wehmütig lief ihm eine rhetorische Frage über die Wange, dann setzte es seinen Weg fort.

Der große Oxmox riet ihr davon ab, da es dort wimmele von bösen Kommata, wilden Fragezeichen und hinterhältigen Semikoli, doch das Blindtextchen ließ sich nicht beirren. Es packte seine sieben Versalien, schob sich sein Initial in den Gürtel und machte sich auf den Weg. Als es die ersten Hügel des Kursivgebirges erklommen hatte, warf es einen letzten Blick zurück auf die Skyline seiner Heimatstadt Buchstabhausen, die Headline von Alphabetdorf und die Subline seiner eigenen Straße, der Zeilengasse. Wehmütig lief ihm eine rhetorische Frage über die Wange, dann setzte es seinen Weg fort.


%%======================================================================================================
%\chapter{Bildeinbindung}
%\label{sec:Bilder}
%
%Weit hinten, hinter den \textbf{Wortbergen}, fern der Länder \textit{Vokalien} und Konsonantien leben die Blindtexte. Abgeschieden wohnen Sie in Buchstabhausen an der Küste\footnote{Landschaft am Meer}
%des Semantik, eines großen \textsl{Sprachozeans}. %Ein kleines Bächlein namens Duden fließt durch ihren Ort und versorgt sie mit den nötigen Regelialien (siehe \ref{sec:Tabellen} S. \pageref{sec:Tabellen}).
%
%\begin{figure}[htbp]
%		\centering
%		\alt{Ich bin das Logo der Technischen Universität Dresden}
%		\includegraphics[scale=0.50]{images/tu_logo}%MCID 8
%	\caption{tu-logo}%MCID 9
%	\label{fig:tulogo}
%\end{figure}
%
%Es ist ein paradiesmatisches Land, in dem einem \texttt{gebratene} Satzteile in den Mund fliegen. Nicht einmal von der allmächtigen Interpunktion\footnote{Zeichensetzung im Satz} werden die Blindtexte beherrscht - ein geradezu unorthographisches Leben. Eines Tages aber beschloss eine \textsc{kleine} Zeile Blindtext, ihr Name war Lorem Ipsum, hinaus zu gehen in die weite Grammatik.
%
%\begin{figure}[htbp]
%		\centering
%		\alt{Ich bin das Logo der Technischen Universität Dresden}
%		\includegraphics[scale=0.10]{images/tu_logo}%MCID 15
%	\caption[kurz]{ein mini TU-Logo}%MCID 16
%	\label{fig:minitulogo}
%\end{figure}
%
%Der große Oxmox riet ihr davon ab, da es dort wimmele von bösen Kommata, wilden Fragezeichen und hinterhältigen Semikoli, doch das Blindtextchen ließ sich nicht beirren. Es packte seine sieben Versalien, schob sich sein Initial in den Gürtel und machte sich auf den Weg.
%
%\begin{figure}[htbp]
%		\centering
%		\alt{Ich bin das Logo der Technischen Universität Dresden}
%		\includegraphics[scale=0.25]{images/tu_logo}%MCID 18
%	\caption{ein kleines TU-Logo}%MCID 19
%	\label{fig:kleinestulogo}
%\end{figure}
%
%Als es die ersten Hügel des Kursivgebirges erklommen hatte, warf es einen letzten Blick zurück auf die Skyline seiner Heimatstadt Buchstabhausen, die Headline von Alphabetdorf und die Subline seiner eigenen Straße, der Zeilengasse. Wehmütig lief ihm eine rhetorische Frage über die Wange, dann setzte es seinen Weg fort.
%
%%======================================================================================================
%\chapter{Zitatumgebungen}
%\label{sec:Zitate}
%
%Weit hinten, hinter den \textbf{Wortbergen}, fern der Länder \textit{Vokalien} und Konsonantien leben die Blindtexte. Abgeschieden wohnen Sie in Buchstabhausen an der Küste\footnote{Landschaft am Meer}
%des Semantik, eines großen \textsl{Sprachozeans}. Ein kleines Bächlein namens Duden fließt durch ihren Ort und versorgt sie mit den nötigen Regelialien (siehe \ref{sec:Tabellen} S. \pageref{sec:Tabellen}).
%
%\begin{quotation}
%	"`Ich bin ein Zitat."'
%\end{quotation}
%
%Es ist ein paradiesmatisches Land, in dem einem \texttt{gebratene} Satzteile in den Mund fliegen. Nicht einmal von der allmächtigen Interpunktion\footnote{Zeichensetzung im Satz} werden die Blindtexte beherrscht - ein geradezu unorthographisches Leben. Eines Tages aber beschloss eine \textsc{kleine} Zeile Blindtext, ihr Name war Lorem Ipsum, hinaus zu gehen in die weite Grammatik. \begin{quote}"`Ich bin ein kurzes Zitat."'\end{quote} Der große Oxmox riet ihr davon ab, da es dort wimmele von bösen Kommata, wilden Fragezeichen und hinterhältigen Semikoli, doch das Blindtextchen ließ sich nicht beirren. Es packte seine sieben Versalien, schob sich sein Initial in den Gürtel und machte sich auf den Weg. Als es die ersten Hügel des Kursivgebirges erklommen hatte, warf es einen letzten Blick zurück auf die Skyline seiner Heimatstadt Buchstabhausen, die Headline von Alphabetdorf und die Subline seiner eigenen Straße, der Zeilengasse. Wehmütig lief ihm eine rhetorische Frage über die Wange, dann setzte es seinen Weg fort.
%
%\begin{verse}
%	Ich bin ein Gedicht\\
%	zumindest mag ich eins werden\\
%	wenn ich mal groß bin.
%\end{verse}
%
%Die Copy warnte das Blindtextchen, da, wo sie herkäme wäre sie zigmal umgeschrieben worden und alles, was von ihrem Ursprung noch übrig wäre, sei das Wort "`und"'  und das Blindtextchen solle umkehren und wieder in sein eigenes, sicheres Land zurückkehren.
%
%%======================================================================================================
%\chapter{Tabellen}
%\label{sec:Tabellen}
%
%Weit hinten, hinter den \textbf{Wortbergen}, fern der Länder \textit{Vokalien} und Konsonantien leben die Blindtexte. Abgeschieden wohnen Sie in Buchstabhausen an der Küste\footnote{Landschaft am Meer}
%des Semantik, eines großen \textsl{Sprachozeans}. Ein kleines Bächlein namens Duden fließt durch ihren Ort und versorgt sie mit den nötigen Regelialien (siehe \ref{sec:Tabellen} S. \pageref{sec:Tabellen}).
%
%\begin{table}[htbp]
%	\centering
%		\begin{tabular}{l|l l}
%			\thead{11} & \thead{12} & \thead{13} \\ \hline
%			21 & 22 & 23 \\
%			31 & 32 & 33
%		\end{tabular}
%	\caption{meine ungetaggte Tabelle}
%	\label{tab:meineUngetaggteTabelle}
%\end{table}
%
%Es ist ein paradiesmatisches Land, in dem einem \texttt{gebratene} Satzteile in den Mund fliegen. Nicht einmal von der allmächtigen Interpunktion\footnote{Zeichensetzung im Satz} werden die Blindtexte beherrscht - ein geradezu unorthographisches Leben. Eines Tages aber beschloss eine \textsc{kleine} Zeile Blindtext, ihr Name war Lorem Ipsum, hinaus zu gehen in die weite Grammatik.
%
%\begin{table}[htbp]
%	\centering
%		\begin{tabular}{l l l}
%			%\toprule
%			\thead{11} & 12 & 13 \\ %\midrule
%			\thead{21} & 22 & 23 \\
%%			\cmidrule(r){1-2}
%			\thead{31} & 32 & 33
%%			\bottomrule
%		\end{tabular}
%	\caption{meine booktabs Tabelle}
%	\label{tab:meinebooktabsTabelle}
%\end{table}
%
%
%Der große Oxmox riet ihr davon ab, da es dort wimmele von bösen Kommata, wilden Fragezeichen und hinterhältigen Semikoli, doch das Blindtextchen ließ sich nicht beirren. Es packte seine sieben Versalien, schob sich sein Initial in den Gürtel und machte sich auf den Weg.
%
\end{document}
